\documentclass[10pt,a4paper]{amsart}
\usepackage[margin=19mm]{geometry}
\usepackage{amsmath,amssymb,mathtools}
\usepackage[scr=rsfs,cal=euler]{mathalfa}
\usepackage{xfrac}
\usepackage[unicode,hidelinks,bookmarks=false]{hyperref}
\newcommand{\ie}{i.e.\ }
\newcommand{\cf}{cf.\ }
\newcommand{\resp}{respectively\ }
\newcommand{\Subsection}{\subsection}
\providecommand{\nobreakdash}{\nobreak\hbox{-}}
\setlength{\emergencystretch}{2em}
\multlinegap0pt
\usepackage{bm}
\usepackage{bbm}
\usepackage{dsfont}
\usepackage{xspace}
\usepackage{cancel}
\usepackage{mathtools}
\usepackage{amsthm}
\makeatletter
\@ifundefined{theorem}{\newtheorem{theorem}{Theorem}}{}
\@ifundefined{assumption}{\newtheorem{assumption}[theorem]{Assumption}}{}
\@ifundefined{definition}{\newtheorem{definition}[theorem]{Definition}}{}
\@ifundefined{proposition}{\newtheorem{proposition}[theorem]{Proposition}}{}
\makeatother
\DeclareMathOperator{\supp}{supp} 
\def\cS{\mathcal{S}}
\title[Non-conservative optimal transport]{Non-conservative optimal transport}
\author{Gabriela Kov\'a\v{c}ov\'a, Georg Menz, Niket Patel}
\date{Source: arXiv:2510.03332v1 (CC BY 4.0)}
\begin{document}
\maketitle

\noindent\emph{Source and licence.} Non-conservative optimal transport. Version arXiv:2510.03332v1 (CC BY 4.0), distributed under the Creative Commons Attribution 4.0 International licence, as recorded in the arXiv licence metadata for this exact version and shown on the abstract page https://arxiv.org/abs/2510.03332v1. Full rights notice: licenses/arxiv-2510.03332-CC-BY-4.0.txt.\par\smallskip

\noindent\textit{Editorial scope.} Counted unit: the Proposition whose source label is \texttt{the:existence\_optimal\_maps} in the article (source0.tex lines 1018--1040, source label \texttt{the:existence\_optimal\_maps}), delivered together with the article's own complete original proof of it (source0.tex lines 1042--1103, one \texttt{proof} environment of 62 lines). No mathematics is written, repaired or re-derived: statement and proof are the article's own text line for line. Required context retained uncounted: lossyOTplans (lines 594--599); equ:katorovich\_problem (lines 596--598); the:solution\_KP (lines 608--613); equ:dual\_KP (lines 715--718); asu:existence\_optimal\_potentials (lines 744--746); pro:solution\_dual\_KP (lines 748--761); the:strong\_duality (lines 895--897); the:characterization\_support\_optimal\_transport\_plan (lines 975--998); equ:criterion\_optimal\_transport\_map (lines 986--988); equ:criterion\_dual\_optimal\_transport\_map (lines 990--992); equ:support\_pi\_characterized via minima (lines 994--996). Every definition, display and label of those lines is retained unchanged. Omitted from this package and not redistributed: 1--593, 599--607, 614--714, 719--743, 747--747, 762--894, 898--974, 989--989, 993--993, 997--1017, 1041--1041, 1104--1676. Nothing inside a retained excerpt is omitted or altered: every retained body is delivered exactly as the article's lines are written, so no omission receipt of any kind accompanies this package. Citations: the retained excerpts carry no citation command at all, so no bibliography entry of the article is transcribed and no reference is redistributed. Reference adaptation: none was needed and none was made; every reference inside the delivered unit resolves inside it, so no unresolved reference or citation is produced. Printed slips kept literal: no printed word, symbol, hypothesis or number of the counted statement or of its proof was altered. Layout and class: the article's own class is \texttt{article} with packages including amsthm, amsmath, amssymb, amsfonts, mathtools, dsfont, enumitem, booktabs, tikz, pgf-pie, comment, subcaption, graphicx, xcolor; the wrapper is the lane's pinned \texttt{amsart} editorial wrapper at 10pt on A4, which restates the theorem-like environments the retained text needs and adds the article's own macro definitions that the retained text needs (\texttt{\textbackslash cS}, \texttt{\textbackslash supp}), with extra wrapper packages (bm, bbm, dsfont, xspace, cancel, mathtools). The delivered numbering is the unit's own. Assets: no retained span contains a figure environment, a graphics-inclusion command, a PDF-page inclusion, a table or a TikZ picture; no image file is copied into this package, the package declares no asset dependency and no image file is redistributed with it. Licence: version arXiv:2510.03332v1 of this article is distributed under the Creative Commons Attribution 4.0 International licence, as recorded in the arXiv licence metadata for this exact version and shown on the abstract page https://arxiv.org/abs/2510.03332v1. A full rights notice is delivered at licenses/arxiv-2510.03332-CC-BY-4.0.txt. Characters: every retained excerpt is pure ASCII, so the delivered engine, XeLaTeX with its default Latin Modern text font, reports no missing character.
\begin{definition}[Kantorovich problem in non-conservative optimal transport] \label{lossyOTplans}
 Given a cost function~$c: \mathcal X \times \mathcal{Y} \to \mathbb{R}$ and  a mass-change factor~$m: \mathcal{X} \times \mathcal{Y} \to \mathbb{R}$, optimize the expression
\begin{align}
\inf_{\pi \in \Gamma_m(\mu, \nu)}  \ \int_{\mathcal{X} \times \mathcal{Y} } c(x,y) \pi(dx,dy) .   \tag{KP}   \label{equ:katorovich_problem}
\end{align}
\end{definition}

\begin{theorem}[Solution of~\eqref{equ:katorovich_problem}]\label{the:solution_KP}
	Suppose that the cost function~$c$ is lower semi-continuous and bounded from below, and that the mass-change factor~$m$ is continuous and bounded. Then the Problem \ref{lossyOTplans} admits a minimizer~$\pi^* \in \Gamma_m(\mu, \nu) $, i.e.
    \begin{align}
        \inf_{\pi \in \Gamma_m(\mu, \nu)}  \ \int_{\mathcal{X} \times \mathcal{Y} } c(x,y) \pi(dx,dy)  =   \int_{\mathcal{X} \times \mathcal{Y} } c(x,y) \pi^*(dx,dy).
    \end{align}
\end{theorem}

\begin{align}
  \sup_{(\varphi, \psi) \in \mathcal{A}(c,m)}&
 \int \varphi(x)\,\mu(dx) . \tag{DP} \label{equ:dual_KP} 
\end{align}

\begin{assumption}\label{asu:existence_optimal_potentials}
  The sets~$\mathcal{X}$ and~$\mathcal{Y}$ are compact and the functions~$c$ and~$m$ are continuous with~$\inf\limits_{x,y}m(x,y)>0$.   
\end{assumption}

\begin{theorem}\label{pro:solution_dual_KP}
 Let Assumption~\ref{asu:existence_optimal_potentials} be satisfied. Then, there exists $(\varphi, \psi)\in \mathcal{A}$ satisfying 
 \begin{align}\label{def:generalized_c_transform}
     \varphi(x) : =   \psi_m^c (x): = \inf_{y \in \mathcal{Y}} \left(c(x,y) - \left( \psi(y) - \int_{\mathcal{Y}} \psi(b) \nu(db) \right) m(x,y) \right)
 \end{align}
 and
 \begin{align}\label{equ:realtion_phi_psi}
        \psi(y) = \inf_{x\in \mathcal{X}} \frac{c(x,y)- \varphi(x)}{m(x,y)} \quad \mbox{with}  \quad \int \psi(y) \nu(dy) = 0.
 \end{align}
 such that
 \begin{align*}
     \sup \eqref{equ:dual_KP} = \max \eqref{equ:dual_KP} = \int \varphi(x)  \mu(dx).
 \end{align*}
\end{theorem}

\begin{theorem}[Strong Duality]\label{the:strong_duality}
If Assumption~\ref{asu:existence_optimal_potentials} holds and $\mathcal X, \mathcal Y$ are convex, then $\inf \eqref{equ:katorovich_problem} = \sup \eqref{equ:dual_KP}$.
\end{theorem}

 \begin{proposition}\label{the:characterization_support_optimal_transport_plan}
          Assume that~$\min \eqref{equ:katorovich_problem} = \max \eqref{equ:dual_KP}$, which means that there is~$\pi^* \in \Gamma_{m} (\mu, \nu)$ and~$(\varphi, \psi) \in \mathcal{A}$ such that
    \begin{align}\label{equ:strong_duality_specific}
       \int \varphi(x) \  \mu(dx) = \int c(x,y)  \ d  \pi^*.
    \end{align}
        Then, the support of~$\pi^*$  is contained in the set
        \begin{align}\label{equ:support_pi_contained_in_c_convex_set}
            \cS :=  \left\{ (x,y) \in \mathcal{X} \times \mathcal{Y} \ | \ c(x,y) = \varphi(x) + \left( \psi(y) - \int \psi \ d \nu \right) m(x,y) \right\},
        \end{align} 
        i.e. $\supp \pi^* \subset \cS$. 
        Moreover, consider the family of functions~$f_x: \mathcal{Y} \to \mathbb{R}$ and~$g_y: \mathcal{X} \to \mathbb{R}$ given by 
        \begin{align}\label{equ:criterion_optimal_transport_map}
            f_x(y):= c(x,y) - \left(\psi(y)- \int \psi \ d \nu \right) m(x,y)
        \end{align}
        and
        \begin{align}\label{equ:criterion_dual_optimal_transport_map}
            g_y(x):= \frac{ c(x,y) - \varphi(x)}{m(x,y)}.
        \end{align}
        Then, the set~$\cS$ is given by
        \begin{align}\label{equ:support_pi_characterized via minima}
            \cS = \left\lbrace (x,y) \in \mathcal{X} \times \mathcal{Y} \ | \ y \mbox{ minimizes } f_x \mbox{ and } x \mbox{ minimizes } g_y \right\rbrace.
        \end{align}
        
    \end{proposition}

\begin{proposition}[Existence of optimal transport maps in the perturbative regime]\label{the:existence_optimal_maps}

    Let $\mathcal{X}, \mathcal{Y} \subset \mathbb{R}^n$ be non-empty, compact, connected sets that are equal to the closure of their interiors (i.e.~$\mathcal{X}, \mathcal{Y}$ are domains). Let~$h,d: [0, \infty) \to \mathbb{R}$ be two smooth, convex, Lipschitz functions such that 
    \begin{itemize}
        \item $h$ is strictly increasing and $\lim_{t \downarrow 0} \frac{h(t) - h(0)}{t} =0$;
        \item $d$ is non-decreasing and $\lim_{t \downarrow 0} \frac{d(t) - d(0)}{t} =0$.
    \end{itemize}
    Additionally, we assume that there is~$\varepsilon>0$
    \begin{align}\label{equ:assumptions_d_and_h}
    \hspace{-0.5cm}
             h''(z)  > \varepsilon d''(z) \quad \mbox{and} \quad
    h'(z) > \varepsilon d'(z)  \qquad \mbox{for all } 0< z \leq \sup_{x\in \mathcal{X} , y \in \mathcal{Y}} |x-y|.      
    \end{align}
    We consider the cost function~$c(x,y):= h(|x-y|)$ and the mass-change factor~$m(x,y)= \exp(- k \ d(|x-y|))$ where $k \geq 0$ is a constant. \newline

    Then, there is a constant~$K$ such that for all $0 \leq k < K$ it holds:
    \begin{enumerate}
    \item[(i)] If the probability measure $\mu \in \mathcal{P} (\mathcal{X})$ absolutely continuous w.r.t.~the Lebesgue measure and the boundary of~$\mathcal{X}$ has Lebesgue-measure 0, then every optimal transport plan is given by a unique optimal transport map~$T: \mathcal{X} \to \mathcal{Y}$.
    
    \item[(ii)] If the probability measure $\nu \in \mathcal{P} (\mathcal{Y})$ absolutely continuous w.r.t.~the Lebesgue measure, and the boundary of~$\mathcal{Y}$ has Lebesgue-measure 0, then then every optimal transport plan is given by a unique dual optimal transport map~$S: \mathcal{Y} \to \mathcal{X}$.
    \end{enumerate}

\end{proposition}

\begin{proof}[Proof of Proposition~\ref{the:existence_optimal_maps}]
Under the current assumptions, Theorems~\ref{the:solution_KP},~\ref{pro:solution_dual_KP}, and~\ref{the:strong_duality} hold. Consequently, the optimal dual solution $(\varphi, \psi) \in \mathcal{A}$ exists and the assumptions of Proposition~\ref{the:characterization_support_optimal_transport_plan} are satisfied. 
Hence, for any point~$(x,y)\in \supp \pi^*$ it holds that the value~$x$ is a minimizer of the function $g_y(x) = \frac{c(x,y)- \varphi(x)}{m(x,y)}$ and the value~$y$ is a minimizer of the function $f_x(y) = c(x,y) -  \psi(y) m(x,y)$. We now show that~$f_x$ and~$g_y$ are a.e.~differentiable.\smallskip

Recall from the proof of Theorem~\ref{pro:solution_dual_KP} that $\varphi, \psi$ are uniformly bounded. Specifically, we can show that 
\begin{align}\label{equ:uniform_bounds_psi_phi}
-\bar{h}(1 + 2 e^{2k\bar{d}}) \leq \varphi (x) \leq \overline{h}  \quad \text{ and } \quad -2 \bar{h} e^{k \bar{d}}  \leq \psi(y) \leq 2 \bar{h} e^{k \bar{d}}    
\end{align} 
where $\overline{h} := \sup\limits_{x \in \mathcal{X}, y \in \mathcal{Y}} | h(|x - y|) |$ and $\overline{d} := \sup\limits_{x \in \mathcal{X}, y \in \mathcal{Y}} |d(|x - y|)|$. In particular those bounds are uniform in~$|k| \leq K$. As~$\mathcal{X}$ and~$\mathcal{Y}$ are compact and additionally~$\varphi$ and~$\psi$ are uniformly bounded, the functions~$\varphi$ and~$\psi$ inherit the modulus of continuity, up to a universal constant, from~$c$ and~$m$ (see proof of Theorem~\ref{pro:solution_dual_KP} for more details). Therefore~$\varphi$ and~$\psi$ are both Lipschitz, which implies together with our assumptions on~$c$ and~$m$ that~$f_x$ and~$g_y$ are also Lipschitz (see~\eqref{equ:criterion_optimal_transport_map} and~\eqref{equ:criterion_dual_optimal_transport_map}). It follows from Rademacher's theorem (here we use that~$\mathcal{X}$ is the closure of its interior) that if~$\mu$ is absolutely continuous w.r.t.~the Lebesgue measure, the functions~$g_y$ are~$\mu$-a.e.~differentiable. Similar the functions~$f_x$ are~$\nu$-a.e.~differentiable if~$\nu$ is absolutely continuous w.r.t.~the Lebesgue measure. \smallskip

Let us turn to the verification of~(i).
Recall from the characterization~\eqref{equ:support_pi_characterized via minima} that for all $(x,y) \in \supp \pi^*$, the value $x$ minimizes the function $g_y$. For~$(x,y) \in \supp \pi^*$ it holds (recall that boundary points are negligible by assumption)
\begin{align}\label{equ:characteristic_equation_optimal_map}
    0 &= m(x,y) \nabla_x g_y(x)  \\ & =  \nabla_x c(x,y) - \nabla_x \varphi (x) - \Big(  c(x,y) - \varphi (x)  \Big)  \nabla_x   \ln m(x,y) \qquad \qquad \mbox{$\mu$ a.e.~$x$.} 
\end{align}
For~$c(x,y)= h(|x-y|)$ and~$m(x,y)= \exp (-k d(|x-y|)) >0$ we obtain
\begin{align}\label{equ:characteristic_equation_1}
    \Big(  h'(|x-y|)   + k  d'(|x-y|) \left(h(|x-y|) - \varphi(x) \right) \Big) \frac{x-y}{|x-y|} =  \nabla_x \varphi (x).
\end{align}

We now argue that after keeping~$x$ fix the equation~\eqref{equ:characteristic_equation_1} uniquely characterizes~$y$. 
For domains $\mathcal{X}, \mathcal{Y}$ there exists a finite $D:= \sup_{x,y} |x-y|$ and since~$(x,y)$ solves the equation~\eqref{equ:characteristic_equation_1}, it must hold that~$x-y = z \ \frac{\nabla_x \varphi(x)}{|\nabla_x \varphi(x)|}$ for some value $-D \leq z \leq D$. In light of~\eqref{equ:characteristic_equation_1}, the scalar $z$ must satisfy
\begin{align}\label{eq:T(z)}
    \frac{z}{|z|} \underbrace{\Big(  h'(|z|) +  k  d'(|z|) \left(h(|z|) - \varphi(x) \right)  \Big)}_{=:F(|z|)} = |\nabla_x \varphi(x)|.
\end{align}
In view of the uniform upper bound $\varphi(x) \leq \bar{h}$ and $z \in [-D, D]$ it holds $h(|z|) - \varphi(x)  \geq -2\bar{h}$. Then, for $0 < k < \frac{\varepsilon}{2 \bar{h}}$ it holds for $|z| > 0$
\begin{align*}
    F(|z|) =  h'(|z|) +  k \  d'(|z|) \left(h(|z|) - \varphi(x) \right)  \geq  h'(|z|) - 2 \bar{h} \ k \  d'(|z|) \overset{\eqref{equ:assumptions_d_and_h}}{>} 0.
\end{align*}
Additionally, it follows form our assumptions that~$F(|z|) = 0 $ if and only if~$|z|=0$. Therefore, in the case~$|\nabla_x \varphi(x)| =0$, it follows that~$y=x$ is uniquely characterized by~$x$. \smallskip

In the case~$|\nabla_x \varphi(x)| \neq 0$, the equation~\eqref{eq:T(z)} may only have solutions $z > 0$ for sufficiently small $k >0$ . We  now show that $F(z)$ is strictly increasing for sufficiently small $k >0$. In view of the assumptions, it holds that
\begin{align}
    F'(z) & = h''(z) +  k   d''(z) \left(h(z) - \varphi(x) \right) + k d'(z) h'(z) \\
    & \geq  h''(z) - 2 \bar{h} k d''(z) > (\varepsilon - 2 \bar{h} k) d''(z) \geq 0,
\end{align} 
for any~$0\leq k\leq \frac{\varepsilon}{2 \bar{h}}$.
This concludes the argument for $y$ being uniquely characterized via~\eqref{eq:T(z)} given $x$, which proves part (i).






The argument for~(ii) follows the footstep of~(i). We have seen above that the functions~$f_x$ are~$\nu$-a.e.~differentiable.  Taking the gradient of~$f_x$ wrt.~$y$  yields for~$(x,y) \in \supp \pi^*$ (boundary points are again negligible by assumption)
\begin{align}\label{equ:characteristic_equation_dual_optimal_map}
    \nabla_y c(x,y) - m(x,y) \ \nabla_y \psi (y)   - \psi (y)  \nabla_y  m(x,y) = 0 \qquad \qquad \mbox{$\nu$ a.e.~$y$.}
\end{align}
    Plugging in the definition of~$c$ and~$m$ and rearranging terms yields
    \begin{align}
        \frac{x-y}{|x-y|} \left( h'(|x-y|) e^{k \ d(|x-y|)} + k \psi(y)  d'(|x-y|)  \right) = - \nabla_y \psi(y).
    \end{align}
    Using the Ansatz~$x-y = z \frac{\nabla_y \psi(y)}{|\nabla_y \psi(y)|}$ yields
\[
\frac{z}{|z|}  \left(  h'(|z|  ) e^{k d(|z| )}   + k  \psi (y)    d'(|z|)  \right) = - | \nabla_y \psi(y)|.
\]
    Due to~\eqref{equ:uniform_bounds_psi_phi}, the function~$|\psi|$ is uniformly bounded independent of~$k$. Similar to the argument in part (i) one can show that for~$0<k$ small enough
    \begin{align*}
       \left(  h'(|z|  ) e^{k d(|z| )}   + k  \psi (y)    d'(|z|)  \right) > 0~ \qquad \mbox{ if }  |z| >0,
    \end{align*}
    and the let hand side can only be zero at~$|z|=0$. This means in the case~$| \nabla_y \psi(y)| =0$, there is only one valid choice~$x=y$. In the case~$| \nabla_y \psi(y)| \neq 0$, the solutions require~$z < 0$. Now, one can argue that for fixed~$y$ the function in the bracket is strictly monotone increasing in~$|z|$ provided~$k$ small enough and therefore has at most one solution. This characterizes~$x$ uniquely in terms of~$y$ and closes the argument.  
\end{proof}

\end{document}
